\documentclass[12pt,a4paper,oneside]{scrreprt}

% Tabellenfarben frueh aktivieren, bevor xcolor geladen wird.
\PassOptionsToPackage{table}{xcolor}

% Zentrale Stil-Datei. Allgemeines Layout und PM-Geruestbausteine liegen gemeinsam in settings.sty.
\usepackage{settings}

% Abbildungs- und Tabellenverzeichnis koennen bei Bedarf zugeschaltet werden.
\KOMAoptions{toc=listof}

% ---------------------------------------------------------------------------
% Ordnerstruktur
% ---------------------------------------------------------------------------
\SetGraphicsFolder{graphics}
\renewcommand{\TextFolder}{texte/}

% ---------------------------------------------------------------------------
% Seitenraender
% Reihenfolge: links, rechts, oben, unten.
% Diese vier Werte koennen direkt hier angepasst werden.
% ---------------------------------------------------------------------------
% ---------------------------------------------------------------------------
% Seitenlayout
% ---------------------------------------------------------------------------

% Abstand vom Blattrand
\newcommand{\SeitenrandOben}{3mm}
\newcommand{\SeitenrandUnten}{3mm}
\newcommand{\SeitenrandLinks}{10mm}
\newcommand{\SeitenrandRechts}{10mm}

% Einstellungen auf das gesamte Dokument anwenden
% Reihenfolge des Befehls: links, rechts, oben, unten
\SetPageMargins
{\SeitenrandLinks}
{\SeitenrandRechts}
{\SeitenrandOben}
{\SeitenrandUnten}
% ---------------------------------------------------------------------------
% Dokumentdaten: nur Platzhalter, keine Projektinhalte
% ---------------------------------------------------------------------------
\renewcommand{\DocumentType}{Projektdokumentation}
%\renewcommand{\DocumentTitle}{[Projekttitel]}
%\renewcommand{\DocumentSubtitle}{[Untertitel]}
\renewcommand{\DocumentUnit}{}
\renewcommand{\DocumentTopic}{}
\renewcommand{\DocumentDescription}{}

\renewcommand{\InstitutionName}{Fakultät für Elektrische Energietechnik}
\renewcommand{\DepartmentName}{[Informatik und Elektrotechnik]}
\renewcommand{\CourseName}{[Modul / Kurs]}
\renewcommand{\TermName}{[WiSe 26/27]}
\renewcommand{\InstructorName}{Prof. Dr. Patrick Henning}
\renewcommand{\SubmissionDate}{\today}
\renewcommand{\GroupName}{}

\renewcommand{\DocumentShortTitle}{[Zweichachsig nachgeführtes PV Modell]}
\renewcommand{\DocumentFooterLabel}{[Projekt]}
\renewcommand{\DocumentLogoPath}{graphics/haw_logo.png}
\renewcommand{\DocumentLogoXOffset}{-20mm} % positiv = rechts, negativ = links
\renewcommand{\DocumentLogoYOffset}{0mm} % positiv = oben, negativ = unten
\renewcommand{\DocumentLogoWidth}{80mm} % Groesse des Logos auf der Titelseite

% ---------------------------------------------------------------------------
% Bearbeiter:innen
% Weitere Zeilen koennen durch Entfernen des Kommentarzeichens aktiviert werden.
% ---------------------------------------------------------------------------
\SetParticipantOne{Jannis Graeser}{947813}{jannis.graeser@stu.haw-kiel.de}
\SetParticipantTwo{Dominic M. Eriksson}{945747}{dominic.m.eriksson@stu.haw-kiel.de}
\SetParticipantThree{}{}{}
\SetParticipantFour{}{}{}
\SetParticipantFive{}{}{}
\SetParticipantSix{}{}{}
\SetParticipantSeven{}{}{}
\SetParticipantEight{}{}{}
\SetParticipantNine{}{}{}
\SetParticipantTen{}{}{}

% Kurze Angaben in Kopf- und Fusszeile vermeiden Ueberlaenge bei grossen Gruppen.
\SetHeader{\ResolvedDocumentShortTitle}{\ResolvedDocumentFooterLabel}{\pagemark}
\SetFooter{\GroupName}{}{Seite \thepage\ von \pageref{LastPage}}

\begin{document}

    \hypersetup{pageanchor=false}
    \MakeDocumentTitle

% --------------------- Vorspann ---------------------
    \pagenumbering{Roman}
    \tableofcontents

% Optional einschalten, sobald Abbildungen oder Tabellen vorhanden sind:
%   \listoffigures
%   \listoftables
%   \listequations
    \clearpage

% --------------------- Hauptteil --------------------
    \StartMainMatter

    \chapter{Einleitung}\label{ch:einleitung}


\section{Ausganssituation und Motivation}\label{sec:ausganssituation-und-motivation}

Zu nachgeführten Photovoltaikanlagen existieren bereits Arbeiten im Umfeld der HAW Hamburg und der TU Wien.
Als bisher bekannte fachliche Anknüpfungspunkte werden eine Bachelorarbeit Simulationen zur Ertragsoptimierung einer einachsig nachgeführten Photovoltaikanlage von Elena Steffens sowie eine Bachelorthesis zur Sonnennachführung von PV-Modulen von Franz Schlecht und eine Diplomarbeit zur Auslegung und Konstruktion einer einachsig nachgeführten Photovoltaik Anlage von Mathias Klotz genannt.
Diese Arbeiten zeigen, dass das Thema der Nachführung von Photovoltaikanlagen bereits wissenschaftlich und praktisch betrachtet wurde.
Für das vorliegende Projekt stehen jedoch nicht die allgemeine Auslegung neuer PV-Nachführsysteme, sondern die Restaurierung, Modernisierung und Erprobung eines bereits vorhandenen Modells im Mittelpunkt.

\section{Aufgabenstellung}\label{sec:aufgabenstellung}

Im Rahmen des Projekts ist die vorhandene PV-Nachführanlage technisch zu untersuchen, zu dokumentieren und soweit möglich wieder in einen funktionsfähigen Zustand zu versetzen.
Dazu sind zunächst der mechanische, elektrische und elektronische Aufbau sowie die bestehenden Schnittstellen durch Reverse Engineering zu erfassen.
Festgestellte Fehler und Schäden sind systematisch einzugrenzen und durch geeignete Reparatur- oder Modernisierungsmaßnahmen zu beheben.

Zusätzlich ist eine mikrocontrollerbasierte Steuerung zu entwickeln, die das PV-Modul automatisch nach der höchsten Lichteinstrahlung ausrichtet und eine manuelle Bedienung ermöglicht.
Die neue Steuerung soll zur vorhandenen elektrischen Schnittstelle kompatibel ausgeführt und separat montiert werden, sodass ein Wechsel zwischen der ursprünglichen und der neuen Steuerung möglich bleibt.
Abschließend sind die Funktionen der Gesamtanlage zu prüfen und die durchgeführten Arbeiten sowie die erzielten Ergebnisse nachvollziehbar zu dokumentieren.

\section{Zielsetzung}\label{sec:zielsetzung}

Ziel der Arbeit ist es, die vorhandene nachgeführte Photovoltaikanlage technisch zu untersuchen, zu dokumentieren und, soweit technisch möglich, zu restaurieren und zu modernisieren.

\begin{itemize}
    \item eine Projektplanung nach GPM-Vorlage,
    \item eine technische Bestandsaufnahme und Reverse-Engineering-Dokumentation
    \item eine soweit technisch mögliche Restaurierung des vorhandenen Modells,
    \item eine separat montierte, pin-kompatible Steuerung für den manuellen und automatischen Betrieb
    \item dokumentierte Funktions-, Bewegungs-, Genauigkeits- und Sicherheitstests,
    \item eine sinnvolle Kurzzeitmessung unter dokumentierten Bedingungen,
    \item eine Dokumentation für Lehr- und Demonstrationszwecke,
    \item eine abschließende Projektdokumentation einschließlich Schaltplänen und Programmcode.
\end{itemize}

\section{Abgrenzung}\label{sec:abgrenzung}

Die Arbeit umfasst sinnvolle Ertragsmessung über mehrere Tage.
Dies beinhaltet keine Simulationen oder Langzeitmessungen.
Die Tests werden bewusst kurzfristig geplant und umgesetzt, da sie abhängig von Wetter, Verfügbarkeit der Anlage und geeigneten Messbedingungen sind.
Außerdem ist nicht vorgesehen, eine vollständig neue kommerzielle PV-Nachführungsanlage zu entwickeln.
Der Schwerpunkt liegt auf der vorhandenen Anlage, ihrer technischen Wiederherstellung und ihrer Nutzbarkeit für Demonstrationszwecke, sowie der ausführlichen Dokumentation, die zusätzlich zu der Projektplanung stattfindet.

\section{Aufbau der Dokumentation}\label{sec:aufbau-der-dokumentation}

Die Dokumentation gliedert sich in einen projektorganisatorischen und einen technischen Teil.
Zunächst werden die Ausgangssituation, die Zielsetzung sowie die Planung und Durchführung des Projekts beschrieben.
Dazu gehören unter anderem die zeitliche Gliederung, die Aufgabenverteilung und der Umgang mit auftretenden Risiken.

Anschließend wird die technische Bearbeitung der PV-Nachführanlage dargestellt.
Diese umfasst die Untersuchung des Ist-Zustands, das Reverse Engineering, die Fehlersuche sowie die durchgeführten Reparatur- und Modernisierungsmaßnahmen.
Darauf aufbauend werden die Entwicklung, der Aufbau und die Funktionsweise der neuen Steuerung erläutert.
Abschließend folgen die Prüfung der Gesamtanlage, die Bewertung der Projektergebnisse und ein Ausblick auf mögliche weiterführende Arbeiten.

\subsection*{Transparenzhinweis zum Einsatz generativer KI}

Im Rahmen dieses Projekts wurde ChatGPT von OpenAI als unterstützendes Werkzeug eingesetzt.
Die Unterstützung umfasste insbesondere die Formulierung und sprachlichen Überarbeitung einzelner Dokumentationsabschnitte sowie zur Erstellung von \LaTeX{}-Beispielen.
Der Umgang mit generativer KI wurde bereits in der Lehrveranstaltung \enquote{Programmieren} im ersten Semester als mögliches Hilfsmittel für die Programmierung und die Bearbeitung derer Aufgaben thematisiert.
Der Einsatz in diesem Projekt erfolgte auf dieser Grundlage und wird an dieser Stelle offengelegt, um den Umfang der KI-Unterstützung transparent von den eigenen Entscheidungen, Bewertungen und praktischen Tätigkeiten abzugrenzen.

Der Einsatz generativer KI erfolgte, weil einige kleinere Nebenaufgabenbereiche ein hohes Maß an Zeit benötigen würden, und nicht viel zum Projekt beitragen würden.
Die KI-generierten Elemente werden an den jeweiligen Stellen gekennzeichnet und klar abgeschränkt.

Die Verantwortung für die Auswahl, Prüfung und Verwendung der übernommenen Inhalte sowie für die daraus gezogenen Schlussfolgerungen liegt bei der Projektgruppe.
Nicht ungeprüft KI-generierte Ausgaben werden nicht als fachlich korrekt vorausgesetzt oder als Fakten angegeben.



    \chapter{Technischer Ist-Zustand}\label{ch:technischer-ist-zustand}

\section{Vorbereitung}\label{sec:vorbereitung}
Zur Vorbereitung der Erfassung des Technischen-ist-zustandes wurden überlegungen angestellt, wie sich der Prozess des Zuordnens einer ID, sowie relevanten Informationen sowie einem Bild, falls nötig, möglichst strukturiert und übersichtlich darstelen lässt.
Ebenfalls relevent hier war es, dass die Daten möglichst leicht in eine \LaTeX-Tabelle überführt werden können.
Hier wurden drei mögliche Konzepte, die zur effizienten aufnahme der Daten in Betracht gezogen wurden.

Für die Excel-Datei, die Web-App und die APK wurde Generative KI eingesetzt.
Diese hat lediglich ein Großteil des Programmcodes, sowie das Hosting der Web-App übernommen.
Ebenfalls wurden Formulierungen angepasst und Texte auf Gramatik und Rechtschreibung überprüft.
Alle bewertenden Elemente, sowie die Begründung wurde durch ein Gruppenmitglied erstellt und nach überarbeitung und umformulierung auf den vorgegebenen Inhalt überprüft.
Alle Entscheidungen wurden ebenfalls aufgrund der eigenen Bewertungsmatrix getroffen.

Die Excel-Datei bietet eine einfache Einrichtung, eine zuverlässige Offline-Nutzung und einen unkomplizierten Datenexport.
Die Bedienung auf einem Tablet ist aufgrund des Touchsystems allerdings sehr unzuverlässig, da man sich schnell mal vertippt, bzw was Falsches auswählen kann.
Außerdem lassen sich komplexe Verbindungen zwischen Bauteilen und zugehörige Fotos nur eingeschränkt übersichtlich erfassen.
Hinzu kommt, dass es für Mobile Endgeräte nur die Office 365 Version von Excel gibt, welche auf Android geräten in ihren Funktionen stark beschränkt ist

Die Web-App ist auf die konkrete Erfassungsaufgabe zugeschnitten und bietet automatische IDs, getrennte Ansichten für Bauteile und Verbindungen sowie große, für Touchscreens geeignete Eingabeelemente.
Ihre Schwächen liegen in der Abhängigkeit vom Browser, der lokalen Speicherung und dem nicht auf jedem Gerät gleichermaßen zuverlässigen Offlinebetrieb.
Ebenfalls zu beachten ist hier, dass man jene erst erstellen muss, welches ein gewisses Grundverständnis von Programmierung mit gezielter nutzung von KI erfordert.

Die APK übernimmt die aufgabenbezogene Bedienoberfläche der Web-App und stellt sie als eigenständige Android-Anwendung bereit.
Sie ermöglicht eine zuverlässige Offline-Nutzung und einen direkten Zugriff auf Gerätefunktionen.
Außerdem kann so eine Umgebung geschaffen werden, die die exakten Bedingungen für das hier geplante vorgehen bereitstellt.
Dem stehen ein höherer Aufwand bei der erstmaligen Erstellung gegenüber.
Aufgrund der hohen Bedienbarkeit und der guten Eignung für die mobile Datenerfassung wurde die APK als bevorzugte Lösung ausgewählt.

Für einen übersichtlichen vergleich werden die Optionen in \cref{tab:nutzwertanalyse-erfassung} aufgelistet und gegeneinnander abgewogen und dann anschließend mit der Formel
\begin{equation}
    N_j =
    \frac{\sum_{i=1}^{n} w_i \cdot b_{ij}}
    {b_{\max} \cdot \sum_{i=1}^{n} w_i}
    \cdot 100\,\%\label{eq:equation}
\end{equation}
bestimmt werden, wobei $w_i$ hier die Prozentuale Gewichtung ist und $b_{ij}$ die selbstgewählte Bewertung, welche von $1 - 5$ gesetzt wurde.

\begin{table}[htbp]
    \centering
    \caption{Nutzwertanalyse der betrachteten Erfassungslösungen}
    \label{tab:nutzwertanalyse-erfassung}
    \begin{tabular}{p{5.3cm}rrrr}
        \hline
        \textbf{Kriterium} &
        \textbf{Gewichtung} &
        \textbf{Excel} &
        \textbf{Web-App} &
        \textbf{APK} \\
        \hline

        Bedienbarkeit auf dem Tablet
        & 20\,\% & 2 & 4 & 5 \\

        Offline-Nutzung
        & 15\,\% & 5 & 3 & 5 \\

        Automatische IDs und Fehlervermeidung
        & 15\,\% & 3 & 5 & 5 \\

        Erfassung von Verbindungen
        & 10\,\% & 2 & 5 & 5 \\

        Einbindung von Fotos
        & 10\,\% & 1 & 3 & 5 \\

        Export für \LaTeX
        & 10\,\% & 5 & 5 & 5 \\

        Einrichtung und Installation
        & 10\,\% & 3 & 4 & 4 \\

        Sicherung und Datenübertragung
        & 5\,\% & 5 & 3 & 4 \\

        Änderungen und Aktualisierungen
        & 5\,\% & 4 & 1 & 1 \\
        \hline

        \textbf{Gesamtbewertung}
        & \textbf{100\,\%}
        & \textbf{63\,\%}
        & \textbf{78\,\%}
        & \textbf{93\,\%} \\

        \hline
    \end{tabular}
\end{table}

mit
\begin{align}
    N_{\text{Excel}}
    &=
    \frac{
        20 \cdot 2 +
        15 \cdot 5 +
        15 \cdot 3 +
        10 \cdot 2 +
        10 \cdot 1 +
        10 \cdot 5 +
        10 \cdot 3 +
        5 \cdot 5 +
        5 \cdot 4
    }{500}
    \cdot 100\,\% \\
    &=
    \frac{315}{500}
    \cdot 100\,\%
    =
    63\,\%,
    \\[1em]
    N_{\text{Web-App}}
    &=
    \frac{
        20 \cdot 4 +
        15 \cdot 3 +
        15 \cdot 5 +
        10 \cdot 5 +
        10 \cdot 3 +
        10 \cdot 5 +
        10 \cdot 4 +
        5 \cdot 3 +
        5 \cdot 1
    }{500}
    \cdot 100\,\% \\
    &=
    \frac{390}{500}
    \cdot 100\,\%
    =
    78\,\%,
    \\[1em]
    N_{\text{APK}}
    &=
    \frac{
        20 \cdot 5 +
        15 \cdot 5 +
        15 \cdot 5 +
        10 \cdot 5 +
        10 \cdot 5 +
        10 \cdot 5 +
        10 \cdot 4 +
        5 \cdot 4 +
        5 \cdot 1
    }{500}
    \cdot 100\,\% \\
    &=
    \frac{465}{500}
    \cdot 100\,\%
    =
    93\,\%.
\end{align}


Im Folgenden werden die Bewertungen der betrachteten Lösungen anhand der zuvor festgelegten Kriterien begründet.


Die Bedienbarkeit erhielt die höchste Gewichtung, da sie einen unmittelbaren Einfluss auf die Geschwindigkeit und Fehleranfälligkeit der Datenerfassung hat.
Bei einer großen Anzahl von Bauteilen können umständliche oder häufig wiederkehrende Eingabeschritte den Arbeitsablauf erheblich verlangsamen.

Die Offline-Nutzung und der Automatisierungsgrad wurden ebenfalls als wichtige Kriterien berücksichtigt.
Eine zuverlässige Offline-Funktion ermöglicht die Erfassung unabhängig von der Qualität der verfügbaren Internetverbindung und verringert das Risiko, dass Arbeitsfortschritte aufgrund von Verbindungsproblemen nicht vollständig gespeichert werden.
Da die Aufnahme einer großen Anzahl von Bauteilen viele gleichartige Arbeitsschritte umfasst, können zudem Unaufmerksamkeit und Ermüdung die Fehlerwahrscheinlichkeit erhöhen.
Eine automatische und fortlaufende Vergabe der Identifikationsnummern sowie vorgegebene Eingabe- und Auswahlmöglichkeiten reduzieren dieses Fehlerpotenzial und gewährleisten eine konsistente Datenerfassung.

Die Unterstützung bei der Erfassung von Verbindungen, die Einbindung von Fotos, der Export für \LaTeX{} sowie der Einrichtungsaufwand wurden im Vergleich zu Bedienbarkeit, Offline-Nutzung und Automatisierung geringer gewichtet.
Diese Funktionen erleichtern den Arbeitsablauf, sind für die grundlegende Aufnahme der Bauteile jedoch nicht in jedem Fall zwingend erforderlich und könnten bei Bedarf durch manuelle Arbeitsschritte ersetzt werden.
Da die Erfassung einer großen Anzahl von Elementen viele wiederkehrende Tätigkeiten umfasst, bieten solche Komfortfunktionen dennoch einen deutlichen Mehrwert.
Sie reduzieren den erforderlichen Arbeitsaufwand, vereinfachen die Dokumentation und senken die Wahrscheinlichkeit von Fehlern, die durch Unaufmerksamkeit oder Ermüdung entstehen können.

Die Kriterien „Sicherung und Datenübertragung“ sowie „Änderungen und Aktualisierungen“ erhielten eine vergleichsweise geringe Gewichtung.
Der Export und die Übertragung der erfassten Daten stellen bei allen betrachteten Lösungen grundsätzlich keinen größeren Aufwand dar.
Eine regelmäßige Sicherung ist zwar wichtig, beeinflusst den eigentlichen Vorgang der Katalogisierung jedoch nur indirekt.
Auch die nachträgliche Anpassbarkeit wurde geringer gewichtet, da die Excel-Datei, die Web-App und die Android-App bereits vor ihrem praktischen Einsatz getestet und auf mögliche Verbesserungen überprüft werden.
Dennoch können während der Erfassung unerwartete Probleme oder neue Anforderungen auftreten.
In diesem Fall unterscheiden sich die Lösungen darin, wie schnell Änderungen vorgenommen und vor Ort bereitgestellt werden können.
Insbesondere bei der Android-App ist dafür die Erstellung und Installation einer aktualisierten APK erforderlich.

Bei der Bewertung von Excel ist außerdem zu berücksichtigen, dass es sich um eine herstellerspezifische Software des Unternehmens Microsoft handelt.
Der Quellcode ist nicht öffentlich zugänglich, und die Weiterentwicklung, Lizenzierung sowie langfristige Verfügbarkeit der Anwendung werden durch den Hersteller bestimmt.
Je nach eingesetzter Version können zudem ein Microsoft-Konto, eine kostenpflichtige Lizenz oder die Nutzung von Cloud-Diensten erforderlich sein.
Daraus entstehen mögliche Abhängigkeiten vom Anbieter und Fragestellungen hinsichtlich Datenschutz und Datenhoheit.
Da die Tabellen grundsätzlich lokal gespeichert und in offene Formate wie CSV exportiert werden können, wurde Excel dennoch als mögliche Lösung berücksichtigt.

\newpage
\subsubsection{Excel-Datei}

\textbf{Bedienbarkeit auf dem Tablet (2/5):}
\begin{itemize}[leftmargin=2em, nosep]
    \item[\pluspunkt] Die tabellarische Darstellung bietet eine kompakte
    Übersicht über die erfassten Daten.
    \item[\minuspunkt] Kleine Zellen und Bedienelemente erschweren die
    Bedienung auf einem Touchscreen.
    \item[\minuspunkt] Wiederkehrende Eingaben und der Wechsel zwischen
    Zellen verlangsamen den Arbeitsablauf.
    \item[\minuspunkt]
    Schäbiges App Layout.
\end{itemize}

\medskip
\textbf{Offline-Nutzung (5/5):}
\begin{itemize}[leftmargin=2em, nosep]
    \item[\pluspunkt] Eine lokal gespeicherte Datei kann vollständig ohne
    Internetverbindung bearbeitet werden.
\end{itemize}

\medskip
\textbf{Automatische IDs und Fehlervermeidung (3/5):}
\begin{itemize}[leftmargin=2em, nosep]
    \item[\pluspunkt] IDs können durch Formeln oder vorbereitete
    Zahlenreihen erzeugt werden.
    \item[\minuspunkt] Formeln und IDs können versehentlich verändert,
    gelöscht oder überschrieben werden.
    \item[\minuspunkt]
    In der Excel App sind manche Formeln und Funktionen nicht möglich, was das Automatisieren beschränkt.
\end{itemize}

\medskip
\textbf{Erfassung von Verbindungen (2/5):}
\begin{itemize}[leftmargin=2em, nosep]
    \item[\pluspunkt] Verbindungen können in einem zweiten Tabellenblatt
    dokumentiert werden.
    \item[\minuspunkt] Die Auswahl und Zuordnung der beteiligten Bauteile
    ist auf dem Tablet vergleichsweise umständlich.
    \item[\minuspunkt]
    Man muss zwischen Tabs wechseln, was den Workflow stört.
    \item[\minuspunkt]
    Die IDs sind nicht vorbereitet, man muss alle selber eingeben.
\end{itemize}

\medskip
\textbf{Einbindung von Fotos (1/5):}
\begin{itemize}[leftmargin=2em, nosep]
    \item[\minuspunkt] Fotos lassen sich nur umständlich in Tabellenzellen
    einfügen und den Bauteilen zuordnen.
    \item[\minuspunkt] Eingefügte Bilder können die Dateigröße und
    Übersichtlichkeit deutlich beeinträchtigen.
    \item[\minuspunkt]
    Das Bewegen von Fotos ist aufgrund der Touchfunktion und die automatische skalierung von Excel ebenfalls komplizierter.
\end{itemize}


\medskip
\textbf{Export für \LaTeX{} (5/5):}
\begin{itemize}[leftmargin=2em, nosep]
    \item[\pluspunkt] Tabellenblätter können direkt als CSV-Dateien
    exportiert und in \LaTeX{} eingelesen werden.
\end{itemize}

\medskip
\textbf{Einrichtung und Installation (3/5):}
\begin{itemize}[leftmargin=2em, nosep]
    \item[\pluspunkt] Nach Erstellung der Tabelle kann diese theoretisch unmittelbar
    verwendet werden.
    \item[\minuspunkt] Auswahlfelder, Formatierungen und Prüfregeln müssen
    zunächst eingerichtet werden.
    \item[\minuspunkt] Excel ist ein herstellerspezifische Produkt von Microsoft, welches einen Account benötigt, der aufgrund eines Problems mit der Studi Lizenz nur erschwert zu bekommen ist.
\end{itemize}

\medskip
\textbf{Sicherung und Datenübertragung (5/5):}
\begin{itemize}[leftmargin=2em, nosep]
    \item[\pluspunkt] Die gesamte Erfassung befindet sich in einer einzelnen
    Datei, die einfach kopiert und gesichert werden kann.
\end{itemize}

\medskip
\textbf{Änderungen und Aktualisierungen (4/5):}
\begin{itemize}[leftmargin=2em, nosep]
    \item[\pluspunkt] Spalten, Kategorien und Auswahlfelder können
    nachträglich angepasst werden.
    \item[\minuspunkt] Änderungen an der Tabellenstruktur können bestehende
    Formeln oder Verweise beeinträchtigen.
\end{itemize}

\newpage

\subsubsection{Web-App}

\textbf{Bedienbarkeit auf dem Tablet (4/5):}
\begin{itemize}[leftmargin=2em, nosep]
    \item[\pluspunkt] Große Eingabefelder und Schaltflächen erleichtern die
    Bedienung auf einem Touchscreen.
    \item[\pluspunkt] Bauteile und Verbindungen werden in getrennten
    Ansichten dargestellt.
    \item[\minuspunkt] Die Bedienung erfolgt innerhalb eines Browsers und
    kann durch dessen Benutzeroberfläche beeinflusst werden.
\end{itemize}

\medskip
\textbf{Offline-Nutzung (3/5):}
\begin{itemize}[leftmargin=2em, nosep]
    \item[\pluspunkt] Nach der ersten Einrichtung kann die Anwendung
    grundsätzlich offline verfügbar gemacht werden.
    \item[\minuspunkt] Der Offlinebetrieb hängt vom Browser und dessen
    Zwischenspeicher ab.
\end{itemize}

\medskip
\textbf{Automatische IDs und Fehlervermeidung (5/5):}
\begin{itemize}[leftmargin=2em, nosep]
    \item[\pluspunkt] Eindeutige IDs werden automatisch vergeben und bleiben
    beim Bearbeiten eines Bauteils erhalten.
    \item[\pluspunkt] Auswahlfelder verhindern viele fehlerhafte oder
    uneinheitliche Eingaben.
\end{itemize}

\medskip
\textbf{Erfassung von Verbindungen (5/5):}
\begin{itemize}[leftmargin=2em, nosep]
    \item[\pluspunkt] Bereits erfasste Bauteile können anhand ihrer ID und
    Bezeichnung ausgewählt werden.
    \item[\pluspunkt] Die Anwendung verhindert die Verbindung eines Bauteils
    mit sich selbst.
    \item[\pluspunkt] Die Anwendung bietet bei Verbindungen nur bislang erstellte IDs an.
\end{itemize}

\medskip
\textbf{Einbindung von Fotos (3/5):}
\begin{itemize}[leftmargin=2em, nosep]
    \item[\pluspunkt] Fotos können direkt ausgewählt und einem Bauteil
    zugeordnet werden.
    \item[\minuspunkt] Die Speicherung vieler Fotos kann den verfügbaren
    Browserspeicher stark beanspruchen.
    \item[\minuspunkt] Die UI für die Fotoauswahl im Browser ist nicht top of the shelf.
\end{itemize}

\medskip
\textbf{Export für \LaTeX{} (5/5):}
\begin{itemize}[leftmargin=2em, nosep]
    \item[\pluspunkt] Bauteile und Verbindungen können getrennt als
    CSV-Dateien exportiert werden.
    \item[\pluspunkt] Es kann ein Backup erstellt werden, was das arbeiten auch auf mehrere Tage problemslos aufteilen lässt .
\end{itemize}

\medskip
\textbf{Einrichtung und Installation (4/5):}
\begin{itemize}[leftmargin=2em, nosep]
    \item[\pluspunkt] Die Anwendung Anwendung kann über einen Link geöffnet
    werden und benötigt keine klassische Installation.
    \item[\minuspunkt] Für den ersten Zugriff und die Vorbereitung des
    Offlinebetriebs ist eine Internetverbindung erforderlich.
    \item[\minuspunkt] Es ist ein OpenAI Account notwendig, um auf die Wesite zuzugreifen.
\end{itemize}

\medskip
\textbf{Sicherung und Datenübertragung (3/5):}
\begin{itemize}[leftmargin=2em, nosep]
    \item[\pluspunkt] Die erfassten Daten können als Sicherungsdatei
    heruntergeladen werden.
    \item[\minuspunkt] Die Daten sind zunächst an den jeweiligen Browser
    gebunden und müssen manuell gesichert werden.
\end{itemize}

\medskip
\textbf{Änderungen und Aktualisierungen (1/5):}
\begin{itemize}[leftmargin=2em, nosep]
    \item[\minuspunkt] Die dokumentierende Person kann Aufbau und Funktionen
    der Anwendung nicht selbstständig verändern.
    \item[\minuspunkt] Anpassungen erfordern Änderungen am Programmcode und
    eine erneute Veröffentlichung.
\end{itemize}

\newpage

\subsubsection{Android-APK}

\textbf{Bedienbarkeit auf dem Tablet (5/5):}
\begin{itemize}[leftmargin=2em, nosep]
    \item[\pluspunkt] Die Oberfläche wurde gezielt für die Bedienung auf
    einem Android-Tablet entwickelt.
    \item[\pluspunkt] Große Eingabefelder, Schaltflächen und getrennte
    Ansichten ermöglichen einen schnellen Arbeitsablauf.
\end{itemize}

\medskip
\textbf{Offline-Nutzung (5/5):}
\begin{itemize}[leftmargin=2em, nosep]
    \item[\pluspunkt] Alle benötigten Anwendungsdateien befinden sich
    innerhalb der APK und stehen ohne Internetverbindung zur Verfügung.
\end{itemize}

\medskip
\textbf{Automatische IDs und Fehlervermeidung (5/5):}
\begin{itemize}[leftmargin=2em, nosep]
    \item[\pluspunkt] Die IDs werden automatisch und eindeutig vergeben.
    \item[\pluspunkt] Auswahlfelder und Eingabeprüfungen reduzieren
    fehlerhafte Einträge.
\end{itemize}

\medskip
\textbf{Erfassung von Verbindungen (5/5):}
\begin{itemize}[leftmargin=2em, nosep]
    \item[\pluspunkt] Bauteile können bequem über ihre ID und Bezeichnung
    ausgewählt und miteinander verknüpft werden.
    \item[\pluspunkt] Die Verbindungen werden getrennt von den Bauteildaten
    gespeichert und übersichtlich dargestellt.
\end{itemize}

\medskip
\textbf{Einbindung von Fotos (5/5):}
\begin{itemize}[leftmargin=2em, nosep]
    \item[\pluspunkt] Fotos können direkt über die Funktionen des Tablets
    ausgewählt und einem Bauteil zugeordnet werden.
    \item[\pluspunkt] Die aufgenommenen Bilder werden gemeinsam mit den
    übrigen Daten gesichert.
    \item[\pluspunkt] UI zum Auswählen ist sehr nutzerfreundlich.
\end{itemize}

\medskip
\textbf{Export für \LaTeX{} (5/5):}
\begin{itemize}[leftmargin=2em, nosep]
    \item[\pluspunkt] Bauteile und Verbindungen können als getrennte
    CSV-Dateien exportiert werden.
    \item[\pluspunkt] Der Android-Dialog ermöglicht das Speichern und Teilen
    der erzeugten Dateien.
\end{itemize}

\medskip
\textbf{Einrichtung und Installation (4/5):}
\begin{itemize}[leftmargin=2em, nosep]
    \item[\pluspunkt] Nach einmaliger Installation kann die Anwendung direkt
    über das App-Symbol gestartet werden.
    \item[\minuspunkt] Da die Anwendung nicht über den Play Store verteilt
    wird, muss die Installation aus einer externen Quelle einmalig erlaubt
    werden.
\end{itemize}

\medskip
\textbf{Sicherung und Datenübertragung (4/5):}
\begin{itemize}[leftmargin=2em, nosep]
    \item[\pluspunkt] Sicherungsdateien können über den Android-Dialog
    gespeichert oder an ein anderes Gerät übertragen werden.
    \item[\minuspunkt] Die Sicherung muss durch die dokumentierende Person
    manuell durchgeführt werden.
\end{itemize}

\medskip
\textbf{Änderungen und Aktualisierungen (1/5):}
\begin{itemize}[leftmargin=2em, nosep]
    \item[\minuspunkt] Änderungen erfordern eine Anpassung des Programmcodes
    und die Erstellung einer neuen APK, was lokal nicht möglich ist.
\end{itemize}
    \chapter{Analyse und Reverse Engineering}\label{ch:analyse-und-reverse-engineering}
    \chapter{Fehleranalyse und Instandsetzung}\label{ch:fehleranalyse-und-instandsetzung}
    \chapter{Entwicklung der Steuerung}\label{ch:entwicklung-der-steuerung}
    \chapter{Gesamtsystem und Inbetriebnahme}\label{ch:gesamtsystem-und-inbetriebnahme}
    \chapter{Funktionsprüfung und Ergebnisse}\label{ch:funktionsprufung-und-ergebnisse}
    \chapter{Projektverlauf}\label{ch:projektverlauf}

\begin{GPMLandscape}

    \begin{GPMTable}
    {L{1.3cm}|L{2.0cm}|L{2.7cm}|L{4.0cm}|C{1.7cm}|Y|Y|Y}
    {Projekttagebuch}
    {tab:projekttagebuch}
    {
        \TableHead{ID} &
        \TableHead{Zeitraum} &
        \TableHead{Bearbeiter} &
        \TableHead{Tätigkeit} &
        \TableHead{Dauer} &
        \TableHead{Ergebnis} &
        \TableHead{Probleme / Abweichungen} &
        \TableHead{Nächste Schritte}
    }
        
        \PTID &
        \texttt{14.07 -- 15.07.2026} &
        \makecell{
            \textsf{Jannis Graeser}\\
            \textsf{Dominic Eriksson}
        } &
        \textsf{Projektinitialisierung} &
        \texttt{10:30 h} &
        \textsf{Ein Exposé.} &
        \textsf{Enige Details waren nicht zufriedenstellend} &
        \textsf{Exposé überarbeiten und einreichen}
        \\
        \hline

        \PTID &
        \texttt{12.08 -- 15.08.2026} &
        \makecell{
            \textsf{Jannis Graeser}\\
            \textsf{Dominic Eriksson}
        } &
        \textsf{Projektinitialisierung} &
        \texttt{3 h} &
        \textsf{Ein Exposé.} &
        \textsf{} &
        \textsf{Beginn Planungsphase}
        \\
        \hline


        \PTID &
        \texttt{08.07 -- 20.08.2026} &
        \makecell{
            \textsf{Jannis Graeser}\\
            \textsf{Dominic Eriksson}
        } &
        \textsf{Beginn der Grobplanungsphase und erstellung der Datein} &
        \texttt{25:32 h} &
        \textsf{Vollständiges GPM Gerüst, mit Ergebnissen aus der Grobplanung} &
        \textsf{} &
        \textsf{Beginn der Planungsphase}
        \\
        \hline

        \PTID &
        \texttt{20.08 -- 31.01.2027} &
        \makecell{
            \textsf{Jannis Graeser}\\
            \textsf{Dominic Eriksson}
        } &
        \textsf{Beginn der Projektdokumentation und erstellung der Datein} &
        \texttt{:14 h} &
        \textsf{Projektdokumentationsgerüst Gerüst, mit ersten Einträgen} &
        \textsf{Es fehlen einige Informationen, die erst später erarbeitet werden können} &
        \textsf{Abgabe}
        \\
        \hline

        \PTID &
        \texttt{22.08.2026} &
        \makecell{
            \textsf{Jannis Graeser}\\
            \textsf{Dominic Eriksson}
        } &
        \textsf{Erstellung und Einbindung eines Repositorys zum gemeinsamen Austausch von Dateien und Quellcode} &
        \texttt{30\,min} &
        \textsf{Repository mit der Projektplanung und der Projektdokumentation eingerichtet} &
        \textsf{Keine Abweichung} &
        \textsf{Pull- und Push-Vorgang durch alle Projektmitglieder testen und anschließend die gemeinsame Bearbeitung aufnehmen}
        \\
        \hline

        \PTID &
        &
        &
        &
        &
        &
        &
        \\

    \end{GPMTable}

\end{GPMLandscape}
    \chapter{Allgemeine Sicherheit}\label{ch:gefahrdungsbeurteilung}

\begin{table}[H]
    \centering
    \renewcommand{\arraystretch}{1.5}

    \begin{tabularx}{\textwidth}{
        |>{\centering\arraybackslash}m{0.27\textwidth}
        |>{\centering\arraybackslash}X
        |>{\centering\arraybackslash}m{0.27\textwidth}|}
        \hline

        \begin{minipage}[c][3.0cm][c]{\linewidth}
            \centering
            \textbf{Bereich:}\\[0.6em]
            Gefahren die aus dem Modell hervorgehen
        \end{minipage}
        &
        \begin{minipage}[c][3.0cm][c]{\linewidth}
            \centering
            \Large\textbf{Gefährdungsbeurteilung}\\
            in Anlehnung an §§ 5,6 ArbSchG
        \end{minipage}
        &
        \begin{minipage}[c][3.0cm][c]{\linewidth}
            \centering
            \textbf{Projekt:}\\[0.6em]
            Zweiachsig nachgeführtes PV-Modell
        \end{minipage}
        \\

        \hline
    \end{tabularx}\label{tab:table}
\end{table}

\begin{GPMTable}
{L{3.2cm}|L{5.0cm}|L{2.2cm}| Y}
{}
{}
{\TableHead{Gefährdung} &
\TableHead{Gefahr} &
\TableHead{persöhnliche Beurteilung vor Maßnahme} &
\TableHead{Gegenmaßnahmen}}


    \multicolumn{4}{|l|}{\textbf{Mechanische Gefährdungen}} \\
    \hline

    &
    Quetschstellen an beweglichen Teilen &
    Akzeptabel &
    Bewegliche Teile vor der Arbeit abschalten.\\
    \hline

    &
    Verletzungen durch scharfe Kanten &
    Inakzeptabel &
    Bauteile auf scharfe Kanten prüfen und diese gegebenenfalls entschärfen.\\
    \hline

    &
    Einzugsgefahr an Motorwellen &
    Inakzeptabel &
    Motoren vor Arbeit von Spannung trennen. \\
    \hline

    &
    Kippen oder Umstürzen der Anlage &
    Akzeptabel &
    Von Tischkanten fernhalten und auf waagerechtem Untergrund stellen. \\
    \hline


    \multicolumn{4}{|l|}{\textbf{Elektrische Gefährdungen}} \\
    \hline

    &
    Berührung spannungsführender Teile &
    Inakzeptabel &
    Anlage vor Arbeiten spannungsfrei schalten und Spannungsfreiheit prüfen. \\
    \hline

    &
    Kurzschluss nach oder bei Arbeiten an der Verdrahtung &
    Inakzeptabel &
    Verdrahtung vor der Inbetriebnahme kontrollieren und geeignete Absicherung verwenden. \\
    \hline

    &
    Beschädigte Leitungen oder Isolierungen &
    Inakzeptabel &
    Regelmäßige Sichtprüfung und Austausch beschädigter Leitungen. \\
    \hline

    &
    Gefahr durch Kondensatoren &
    Inakzeptabel &
    Vor arbeiten Kondensatoren entladen. \\
    \hline

    &
    Überlastungsgefahr &
    Inakzeptabel &
    Sicherungen einbauen. \\
    \hline


    \multicolumn{4}{|l|}{\textbf{Chemische Gefährdungen}} \\
    \hline

    &
    Hautkontakt mit gesundheitsschädlichen Stoffen &
    Inakzeptabel &
    Nach Berührung der Elemente Händewaschen. \\
    \hline

    &
    Einatmung von gesundheitsschädlichen Dämpfe &
    Inakzeptabel &
    Nur in belüfteten Räumen oder mit Abzug arbeiten. \\
    \hline

    \multicolumn{4}{|l|}{\textbf{Brand und Explosionsgefahr}} \\
    \hline

    &
    Zündquellen durch Lichtbögen &
    Inakzeptabel &
    Alle Kontakte vor einschalten Prüfen. \\
    \hline

    &
    Erwärmung von Bauteilen &
    Akzeptabel &
    Sicherungen zur Verhinderung von Überlastung, sowie Bauteile auf die Belastung auslegen. \\
    \hline


    &
    Explosion von Kondensatoren &
    Inakzeptabel &
    Kondensatoren auf korrekte Polung prüfen. \\
    \hline


\end{GPMTable}
\newpage
\begin{UnifiedInfoBox}
    \textbf{Allgemeine Maßnahmen und der Umgang bei Arbeiten an der Anlage}

    \begin{itemize}
        \item Arbeiten an der elektrischen Anlage grundsätzlich nur im spannungsfreien Zustand durchführen.
        \item Die Anlage vor Beginn der Arbeiten gegen unbeabsichtigtes Einschalten sichern.
        \item Bewegliche Teile und mögliche Quetschstellen während der Arbeiten beachten.
        \item Nur zertifizierte und geeignete Werkzeuge verwenden und beschädigte Werkzeuge nicht einsetzen.
        \item Leitungen, Steckverbindungen und Bauteile vor der Inbetriebnahme auf Beschädigungen prüfen.
        \item Den Arbeitsbereich sauber und frei von unnötigen Gegenständen halten.
        \item Änderungen an der Verdrahtung oder Steuerung vor dem Einschalten kontrollieren.
        \item Die Anlage bei ungewöhnlichem Verhalten, Geruch, Geräuschen oder Erwärmung sofort abschalten.
        \item Praktische Arbeiten nur mit mindestens zwei Personen durchführen.
        \item Kein offenes Schuhwerk tragen.
        \item Keine Nahrung in der Arbeitsumgebung aufbewahren.
        \item Nach Abschluss der Arbeiten alles sicher verstauen und zurückräumen.
    \end{itemize}

\end{UnifiedInfoBox}






% --------------------- Anhang -----------------------
    \appendix
    \input{texte/99_anhang}

% --------------------- Literatur --------------------
% Bei Bedarf aktivieren und Quellen in bib.bib eintragen:
%   \printbibliography[heading=litverz]

\end{document}
